\fancychapter{Temporary}
This chapter contains some content that is not yet defined were will it be used on the document.

\section{Tracking speed}
The main purpose of this project is to create a fast response \ac{mppt}, but how fast does it need to be? To answer this question, an analysis based on real data from \ac{tsb} were conducted. The data was collected from a weekend of test with \ac{sr03}. 

As discussed in previus sections, the \ac{mpp} point can change due to changes in irradiance and temperature, but also due to changes on the angle of incidence of the sun. This change in angle is mainly due to the movement of the boat and the ocean waves. To analyze this effect the pitch and yaw of the boat were collected using an AHRS sensor.




\textcolor{red}{Note:} Realisticaly this data is not reliable because it was taken in 2022, with a different boat of the one we are targeting, which can have a different response due to its mass, center of gravity, size, etc. Also, I do not know the conditions of the test, for example if was taken in a dam, river or ocean. Finally, \ac{sr03} was a boat with hydro foils, which can be aqua planing or just floating, which can also change the response of the boat. It is important to take new data with the targeting boat to have a more accurate analysis. 

\textcolor{red}{Descobri que os dados foram tirados na barragem de montargil e por isso acho que esta analise não é valida}